\chapter{Dahlberg Step 1: the preliminary diffeomorphism (Phase D-B)}
\label{chap:dahlberg_step1}

% archon:covers Gluck/DahlbergStep1.lean

This chapter formalises Step~1 of Dahlberg's proof of the plane case
(\cref{thm:dahlberg_converse}): from the mixed-sign four-vertex hypothesis we
construct a preliminary orientation-preserving circle diffeomorphism $\eta$ and
positive constants $0<a<b$ such that the reparametrised curvature
$\kappa\circ\eta$ agrees, up to a small-$L^1$ error $e$, with the \emph{clean
bicircle} $a(1-f)+bf$. The clean bicircle is strictly positive, so all of the
in-tree positive-case machinery (winding, chord-integral simplicity) applies to
it; Phase D-C (\cref{chap:dahlberg_step2}) then transports the conclusion to
$\kappa\circ\eta$ through the $L^1$ smallness of $e$.

Throughout, $\kappa : \mathbb{R}\to\mathbb{R}$ is continuous and $2\pi$-periodic
(the project's encoding of a function on the circle $\mathbf{T}$), and $f$ is the
fixed two-arc indicator below. Note that, in contrast with the positive case,
$\kappa$ is \emph{not} assumed positive: positivity is hypothesised only at the
two maxima.

\section{The fixed two-arc indicator and the clean bicircle}

% SOURCE: [Dahlberg], §3, p. 2134, definition of $f$ (read from references/dahlberg.pdf)
% SOURCE QUOTE: "Let $f : T \to R$ denote the function defined by $f(e^{i\theta}) = 1$
% whenever $|\theta - \frac{\pi}{2}| < \frac{\pi}{4}$ or $|\theta - \frac{3\pi}{2}| < \frac{\pi}{4}$
% and zero elsewhere."
\begin{definition}

\leanok
[Two-arc indicator $f$]
  \label{def:dahlberg_f}
  \lean{Gluck.dahlbergF}
  \textit{Source: Dahlberg, §3, definition of $f$.}
  Let $f : \mathbb{R}\to\mathbb{R}$ be the $2\pi$-periodic function with
  $f(\theta)=1$ when $|\theta-\tfrac{\pi}{2}|<\tfrac{\pi}{4}$ or
  $|\theta-\tfrac{3\pi}{2}|<\tfrac{\pi}{4}$ (i.e.\ on the open arcs
  $(\tfrac{\pi}{4},\tfrac{3\pi}{4})$ and $(\tfrac{5\pi}{4},\tfrac{7\pi}{4})$) and
  $f(\theta)=0$ elsewhere on $[0,2\pi)$.
\end{definition}

\begin{definition}

\leanok
[Clean bicircle $a(1-f)+bf$]
  \label{def:clean_bicircle}
  \lean{Gluck.cleanBicircle}
  \uses{def:dahlberg_f}
  For reals $a,b$ the \emph{clean bicircle} is the $2\pi$-periodic step function
  \[
    k_{a,b}(\theta) \;=\; a\bigl(1-f(\theta)\bigr) + b\,f(\theta)
      \;=\;
    \begin{cases}
      b, & \theta\in(\tfrac{\pi}{4},\tfrac{3\pi}{4})\cup(\tfrac{5\pi}{4},\tfrac{7\pi}{4}),\\
      a, & \text{otherwise on }[0,2\pi).
    \end{cases}
  \]
\end{definition}

% --- Foundational properties of $f$ and $k_{a,b}$ (coverage-debt helpers) -----
\begin{lemma}

[Periodicity and integrability of $f$ and $k_{a,b}$]
  \label{lem:clean_bicircle_basic}
  \lean{Gluck.dahlbergF_periodic, Gluck.dahlbergF_intervalIntegrable, Gluck.cleanBicircle_periodic, Gluck.cleanBicircle_intervalIntegrable, Gluck.cleanBicircle_bounds}
  \uses{def:dahlberg_f, def:clean_bicircle}
  The indicator $f$ is $2\pi$-periodic and interval-integrable on every $[c,d]$;
  consequently the clean bicircle $k_{a,b}=a(1-f)+bf$ is $2\pi$-periodic and
  interval-integrable, and satisfies the two-sided bound $a\le k_{a,b}\le b$ when
  $a\le b$.
\end{lemma}
\begin{proof}
  \uses{def:dahlberg_f, def:clean_bicircle}
  $f$ is the indicator of a $2\pi$-periodic union of arcs, hence $2\pi$-periodic and
  bounded measurable, so interval-integrable; the bound $0\le f\le 1$ gives
  $a\le a(1-f)+bf\le b$. Periodicity and integrability of $k_{a,b}$ follow by linear
  combination. (Lean helpers: \texttt{dahlbergF\_periodic},
  \texttt{dahlbergF\_intervalIntegrable}, \texttt{cleanBicircle\_periodic},
  \texttt{cleanBicircle\_intervalIntegrable}, \texttt{cleanBicircle\_bounds}.)
\end{proof}

\begin{lemma}

[On-period value of $f$ and the period integrals]
  \label{lem:clean_bicircle_integral}
  \lean{Gluck.integral_cleanBicircle, Gluck.dahlbergF_on_period, Gluck.integral_dahlbergF}
  \uses{def:dahlberg_f, def:clean_bicircle}
  On $[0,2\pi]$, $f$ is the indicator of the two base arcs
  $(\tfrac{\pi}{4},\tfrac{3\pi}{4})\cup(\tfrac{5\pi}{4},\tfrac{7\pi}{4})$, so
  $\int_0^{2\pi}f=\pi$ and $\int_0^{2\pi}k_{a,b}=a\!\cdot\!\pi+b\!\cdot\!\pi=(a+b)\pi$.
\end{lemma}
\begin{proof}
  \uses{def:dahlberg_f, def:clean_bicircle}
  Each value-$1$ arc has length $\pi/2$, total $\pi$ (Lean:
  \texttt{dahlbergF\_on\_period}, \texttt{integral\_dahlbergF}); integrate the linear
  combination $a(1-f)+bf$ to get $(a+b)\pi$.
\end{proof}

\begin{lemma}


\leanok
[Uniform bound of a continuous periodic curvature]
  \label{lem:exists_kappa_bound}
  \lean{Gluck.exists_kappa_bound}
  A continuous $2\pi$-periodic $\kappa:\mathbb{R}\to\mathbb{R}$ is uniformly bounded:
  there is $N$ with $|\kappa(\theta)|\le N$ for all $\theta$.
\end{lemma}
\begin{proof}

\leanok
  Compactness of one period $[0,2\pi]$ bounds $\kappa$ there; periodicity
  (via the $\mathrm{toIcoMod}$ representative) extends the bound to all of $\mathbb{R}$.
\end{proof}

\section{Level extraction from the four-vertex condition}

In the positive case the in-tree \cref{lem:abab_values}
(\texttt{exists\_abab\_of\_fourVertex}) takes the levels $a,b$ to be the extreme
values $\max(\kappa(q_1),\kappa(q_2))$ and $\min(\kappa(p_1),\kappa(p_2))$ and
relies on global positivity ($a>0$ because the minima are positive). In the
\emph{mixed-sign} case the minima may be $\le 0$, so that choice fails. We must
instead pick the two common levels $0<a<b$ \emph{interior} to the gap
$\bigl(\max(0,\max(\kappa(q_1),\kappa(q_2))),\ \min(\kappa(p_1),\kappa(p_2))\bigr)$
--- non-empty because positivity is hypothesised at the maxima and the
value-separation guarantees the gap is open --- and attain them on the four
flanks by the intermediate value theorem. This is the only place the
positivity-at-the-maxima hypothesis is used.

% SOURCE: [Dahlberg], §1, p. 2131, Theorem 1.1 hypothesis (read from references/dahlberg.pdf)
% SOURCE QUOTE: "Suppose $\kappa : T \to R$ is not a constant and assume that
% $\kappa$ has at least four critical points ordered counterclockwise
% $p_i \in T$, $i = 1,\dots,4$, such that $p_1,p_3$ are local maxima and
% $p_2,p_4$ are local minima with $\kappa(p_1) > 0, \kappa(p_3) > 0$ and
% $\max(\kappa(p_2),\kappa(p_4)) < \min(\kappa(p_1),\kappa(p_3))$."
\begin{definition}

\leanok
[Mixed-sign four-vertex condition]
  \label{def:mixed_sign_four_vertex}
  \lean{Gluck.MixedSignFourVertex}
  \uses{def:four_vertex_condition}
  \textit{Source: Dahlberg, §1, Theorem 1.1 hypothesis.}
  A continuous, $2\pi$-periodic, non-constant $\kappa:\mathbb{R}\to\mathbb{R}$
  satisfies the \emph{mixed-sign four-vertex condition} if there are points
  $p_1<q_1<p_2<q_2<p_1+2\pi$ in one fundamental period with $p_1,p_2$ local
  maxima, $q_1,q_2$ local minima, the value separation
  $\max(\kappa(q_1),\kappa(q_2))<\min(\kappa(p_1),\kappa(p_2))$, \emph{and}
  positivity at the maxima $0<\min(\kappa(p_1),\kappa(p_2))$. This is the
  in-tree value-separated \cref{def:four_vertex_condition} strengthened only by
  the maxima-positivity clause; the minima $\kappa(q_1),\kappa(q_2)$ may be
  $\le 0$.
\end{definition}

\begin{lemma}

\leanok
[Strictly positive four-vertex curvature is mixed-sign]
  \label{lem:mixed_sign_four_vertex_of_is_curvature_function}
  \lean{Gluck.mixedSignFourVertex_of_isCurvatureFunction}
  \uses{def:is_curvature_function, def:four_vertex_condition, def:mixed_sign_four_vertex}
  Let $\kappa$ be a curvature function (\cref{def:is_curvature_function}: continuous,
  $2\pi$-periodic and \emph{strictly positive}) which is non-constant and satisfies
  the four-vertex condition (\cref{def:four_vertex_condition}). Then $\kappa$
  satisfies the mixed-sign four-vertex condition
  (\cref{def:mixed_sign_four_vertex}). In particular the mixed-sign hypothesis
  subsumes the non-constant strictly-positive case, so the mixed-sign converse
  applies to every non-constant positive curvature function meeting the four-vertex
  condition.
\end{lemma}
\begin{proof}
  \uses{def:is_curvature_function, def:four_vertex_condition, def:mixed_sign_four_vertex}
  \leanok
  Since $\kappa$ is non-constant, the constant branch of the four-vertex condition
  is excluded, so \cref{def:four_vertex_condition} supplies the four
  value-separated alternating extrema $p_1<q_1<p_2<q_2<p_1+2\pi$ with $p_1,p_2$
  local maxima, $q_1,q_2$ local minima, and
  $\max(\kappa(q_1),\kappa(q_2))<\min(\kappa(p_1),\kappa(p_2))$. These data, together
  with continuity and $2\pi$-periodicity, are exactly the requirements of the
  mixed-sign condition except for the extra positivity-at-the-maxima clause
  $0<\min(\kappa(p_1),\kappa(p_2))$. But that clause is automatic: $\kappa$ is
  everywhere strictly positive, so $0<\kappa(p_1)$ and $0<\kappa(p_2)$, whence
  $0<\min(\kappa(p_1),\kappa(p_2))$. Thus $\kappa$ satisfies
  \cref{def:mixed_sign_four_vertex}.
\end{proof}

% SOURCE: [Dahlberg], §3, p. 2135, first construction paragraph (read from references/dahlberg.pdf)
% SOURCE QUOTE: "We need some definitions in order to construct the first
% preliminary diffeomorphism. For $j = 1,\dots,4$, let $\theta_j = (j-1)\frac{\pi}{2}$
% and set $q_j = e^{i\theta_j}$, $A_j = \{e^{i\theta} : |\theta - \theta_j| < \frac{\pi}{4} - \epsilon\}$
% for a small positive number $\epsilon$. By continuity there are points
% $r_j \in T$, $j = 1,\dots,4$, also ordered counterclockwise and positive numbers
% $a,b$ such that $0 < a < b$ and $\kappa(r_1) = \kappa(r_3) = a$,
% $\kappa(r_2) = \kappa(r_4) = b$. Pick intervals $B_j \subset T$, $j = 1,\dots,4$,
% such that $r_j \in B_j$ and $|\kappa(w) - \kappa(r_j)| < \epsilon$ for all
% $w \in B_j$."
% SOURCE: [Dahlberg], §3, p. 2135, ``By continuity there are points $r_j$'' (read from references/dahlberg.pdf)
% SOURCE QUOTE: "By continuity there are points $r_j \in T$, $j = 1,\dots,4$,
% also ordered counterclockwise and positive numbers $a,b$ such that
% $0 < a < b$ and $\kappa(r_1) = \kappa(r_3) = a$, $\kappa(r_2) = \kappa(r_4) = b$."
\begin{lemma}

\leanok
[Mixed-sign level extraction]
  \label{lem:dahlberg_alignment_data}
  \lean{Gluck.exists_alignmentData}
  \uses{def:mixed_sign_four_vertex, lem:ivt_hits}
  \textit{Source: Dahlberg, §3, ``by continuity there are points $r_j$\ldots''.}
  Let $\kappa$ satisfy the mixed-sign four-vertex condition
  (\cref{def:mixed_sign_four_vertex}). Then there exist constants $0<a<b$ and four
  points $r_1<r_2<r_3<r_4<r_1+2\pi$, ordered counterclockwise, with
  \[
    \kappa(r_1)=\kappa(r_3)=a,\qquad \kappa(r_2)=\kappa(r_4)=b.
  \]
\end{lemma}
\begin{proof}
  \uses{lem:ivt_hits, def:mixed_sign_four_vertex}
  \leanok
  Write $m=\max(\kappa(q_1),\kappa(q_2))$ and $M=\min(\kappa(p_1),\kappa(p_2))$;
  by hypothesis $m<M$ and $0<M$. Pick two levels with
  $\max(0,m)<a<b<M$ (possible since the interval $(\max(0,m),M)$ is non-empty:
  $\max(0,m)<M$ holds because $m<M$ and $0<M$). The four flanks are the closed
  arcs $[p_1,q_1]$, $[q_1,p_2]$, $[p_2,q_2]$, $[q_2,p_1+2\pi]$; on each, $\kappa$
  runs monotonically (in value, not assumed monotone --- the IVT only needs the
  endpoint values to straddle the target) between a maximum value $\ge M>b>a$ and
  a minimum value $\le m<a<b$. Apply the intermediate value theorem
  \cref{lem:ivt_hits} four times:
  \begin{itemize}
    \item on the descent $[p_1,q_1]$ pick $r_1$ with $\kappa(r_1)=a$
      (between $\kappa(q_1)\le m<a$ and $\kappa(p_1)\ge M>a$);
    \item on the ascent $[q_1,p_2]$ pick $r_2$ with $\kappa(r_2)=b$;
    \item on the descent $[p_2,q_2]$ pick $r_3$ with $\kappa(r_3)=a$;
    \item on the ascent $[q_2,p_1+2\pi]$ pick $r_4$ with $\kappa(r_4)=b$
      (using $\kappa(p_1+2\pi)=\kappa(p_1)$ by periodicity).
  \end{itemize}
  Then $p_1\le r_1\le q_1\le r_2\le p_2\le r_3\le q_2\le r_4\le p_1+2\pi$ gives the
  ordering $r_1<r_2<r_3<r_4<r_1+2\pi$ (the inequalities are strict because the
  values $a\ne b$ differ across each interior endpoint $q_i,p_i$), and the value
  pattern is $a,b,a,b$ as required. This is the same four-fold IVT skeleton as
  \cref{lem:abab_values}; only the choice of the levels $a,b$ (interior, rather
  than the extreme values) and the use of maxima-positivity differ.
\end{proof}

\section{The preliminary diffeomorphism and the error decomposition}

% SOURCE: [Dahlberg], §3, p. 2135, $\eta$-construction (read from references/dahlberg.pdf)
% SOURCE QUOTE: "Let $\eta$ be any smooth orientation preserving
% $C^\infty$-diffeomorphism of the circle such that $\eta(A_j) \subset B_j$ for
% $j = 1,\dots,4$. It is now easily seen that $\kappa \circ \eta = a(1-f) + bf + e$,
% where $e$ is bounded with $\int_0^{2\pi} |e| dt < C\epsilon$ and $C$ is
% independent of $\epsilon$."
\begin{theorem}

\leanok
[Step 1: preliminary diffeomorphism]
  \label{thm:exists_preliminary_diffeo}
  \lean{Gluck.exists_preliminaryDiffeo}
  % DESIGN NOTE (iter-031): $e$ is interval-integrable, NOT continuous. Since
  % $\kappa\circ\eta$ is continuous (C¹ $\eta$) and $a(1-f)+bf$ is a step function,
  % $e=\kappa\circ\eta-(a(1-f)+bf)$ has jump discontinuities at the four breakpoints;
  % only its $L^1$ smallness is used downstream. The Lean signature accordingly
  % carries `IntervalIntegrable e volume 0 (2π)` in place of `Continuous e`.
  \uses{def:dahlberg_f, def:clean_bicircle, lem:dahlberg_alignment_data, def:mixed_sign_four_vertex}
  \textit{Source: Dahlberg, §3, Step 1.}
  Let $\kappa$ be continuous, $2\pi$-periodic, non-constant and satisfy the
  mixed-sign four-vertex condition (\cref{def:mixed_sign_four_vertex}). There are
  constants $0<a<b$ (the levels
  of \cref{lem:dahlberg_alignment_data}) and, for every $\varepsilon>0$, a $C^1$
  (in fact $C^\infty$) orientation-preserving circle diffeomorphism $\eta$
  ($\eta(t+2\pi)=\eta(t)+2\pi$, $\eta'>0$) and an interval-integrable,
  $2\pi$-periodic error $e$ with
  \[
    \kappa\circ\eta \;=\; a(1-f) + b\,f + e,
    \qquad \int_0^{2\pi} |e(t)|\,dt \;<\; C\,\varepsilon,
  \]
  where the constant $C$ does not depend on $\varepsilon$. Consequently the total
  curvature $I=\int_0^{2\pi}\kappa\circ\eta$ satisfies
  $I = (a+b)\pi + \int_0^{2\pi}e$, so $|I-(a+b)\pi|<C\varepsilon$ and in
  particular $I>0$ once $\varepsilon$ is small (since $a,b>0$).
\end{theorem}
\begin{proof}
  \uses{lem:dahlberg_alignment_data, lem:step_curvature_periodic, lem:exists_eta_clean_L1, lem:clean_bicircle_basic, lem:clean_bicircle_integral}
  \leanok
  Take $a,b$ and the points $r_j$ from \cref{lem:dahlberg_alignment_data}.
  Following Dahlberg, fix the model breakpoints $\theta_j=(j-1)\tfrac{\pi}{2}$ and
  the shrunken arcs $A_j=\{\theta:|\theta-\theta_j|<\tfrac{\pi}{4}-\varepsilon\}$.
  By continuity of $\kappa$ choose closed arcs $B_j\ni r_j$ with
  $|\kappa(w)-\kappa(r_j)|<\varepsilon$ for $w\in B_j$. Let $\eta$ be any smooth
  orientation-preserving circle diffeomorphism carrying $A_j$ into $B_j$ for
  $j=1,\dots,4$ (such an $\eta$ exists: prescribe $\eta$ on the four model arcs to
  land in $B_j$ and interpolate monotonically and smoothly on the four
  complementary transition gaps, each of length $2\varepsilon$).

  On the bulk $A_j$ the composite $\kappa\circ\eta$ is within $\varepsilon$ of the
  constant $\kappa(r_j)$, which equals the clean-bicircle value $a$ (for
  $j=1,3$) or $b$ (for $j=2,4$). Hence the error $e=\kappa\circ\eta-(a(1-f)+bf)$
  satisfies $|e|<\varepsilon$ on $\bigcup_j A_j$, a set of total measure
  $2\pi-8\varepsilon$, contributing less than $2\pi\varepsilon$ to
  $\int|e|$. On the four transition gaps (total measure $8\varepsilon$) the error
  is bounded by $2M$, where $M=\sup|\kappa|+\max(a,b)$ is a fixed bound
  independent of $\varepsilon$; this contributes less than $16M\varepsilon$. Thus
  $\int_0^{2\pi}|e|<(2\pi+16M)\varepsilon=:C\varepsilon$ with $C$ independent of
  $\varepsilon$.

  Finally $\int_0^{2\pi}(a(1-f)+bf)=a\cdot(2\pi-2\cdot\tfrac{\pi}{2})+b\cdot(2\cdot\tfrac{\pi}{2})
  =a\pi+b\pi=(a+b)\pi$ (each value-$b$ arc has length $\tfrac{\pi}{2}$, two of
  them), so $I=(a+b)\pi+\int e$ and the stated bound on $I$ follows from
  $|\int e|\le\int|e|<C\varepsilon$.
\end{proof}

\begin{lemma}


\leanok
[Analytic core: the $L^1$ $\eta$-construction]
  \label{lem:exists_eta_clean_L1}
  \lean{Gluck.exists_eta_clean_L1}
  \uses{lem:dahlberg_alignment_data, def:clean_bicircle, lem:clean_bicircle_integral, lem:exists_kappa_bound}
  The genuinely-new analytic content of \cref{thm:exists_preliminary_diffeo},
  isolated: for $\kappa$ continuous $2\pi$-periodic and the levels/crossing points
  $r_1<r_2<r_3<r_4$ of \cref{lem:dahlberg_alignment_data}, and every
  $\varepsilon>0$, there is a $C^1$ orientation-preserving quasiperiodic $\eta$
  (explicit continuous positive derivative) with
  $\int_0^{2\pi}|\kappa(\eta\,t)-k_{a,b}(t)|\,dt<\varepsilon$.
  \cref{thm:exists_preliminary_diffeo} is a routine wrapper around this lemma.
\end{lemma}
\begin{proof}
  \uses{lem:dahlberg_alignment_data, lem:exists_kappa_bound}
  \leanok
  Take $\eta(\theta)=h_1(\theta+\tfrac{\pi}{4})$ with $h_1(\theta)=m_0+\int_0^\theta w$,
  $w$ the plateau density of $\mathrm{exists\_plateau\_density}$ at the crossing
  points $r_j$; the $+\tfrac{\pi}{4}$ shift aligns the four model plateaus (centred
  at $0,\tfrac{\pi}{2},\pi,\tfrac{3\pi}{2}$) with the constancy arcs of $k_{a,b}$.
  $\eta$ is $C^1$ with $v=w(\cdot+\tfrac{\pi}{4})>0$ (chain rule) and quasiperiodic
  ($\eta(t+2\pi)=\eta(t)+2\pi$). Recalibrating
  $\delta=\min(\varepsilon'/8,\pi/4,\varepsilon/(8(N{+}b{+}1)))$ with $N$ the uniform
  $\kappa$-bound (\cref{lem:exists_kappa_bound}), the $L^1$ estimate splits $[0,2\pi]$
  into five plateau pieces (integrand $\le\varepsilon'$, clean value via the on-period
  form of $k_{a,b}$) and four transition gaps (integrand $\le N+b$, total gap measure
  $<\varepsilon/2$); combining via $\mathrm{integral\_add\_adjacent\_intervals}$ and
  $\mathrm{integral\_mono\_on}$ gives $\int_0^{2\pi}|\kappa\circ\eta-k_{a,b}|<\varepsilon$.
\end{proof}

\noindent\textbf{Phase status.} \cref{thm:exists_preliminary_diffeo} delivers
exactly the hypothesis that Phase D-C consumes: a reparametrisation $\eta$ after
which the curvature is a strictly positive clean bicircle plus an $L^1$-small
perturbation, with positive total curvature $I>0$. The remaining work --- finding
the \emph{closing} adjustment over the in-tree configuration disk and transporting
closure and simplicity through $\int|e|<C\varepsilon$ --- is
\cref{chap:dahlberg_step2}.
